\chapter{Discussion}
\label{sec:discussion}

\section{The sim-to-real gap}

The sim-to-real drop is measured between matched configurations:
$\SimToRealGapStatic{}$ macro-F1 for the deployed LC+static model
($\TierLCStaticMacroFFull{} \to \WithStaticMacroF{}$) and
$\SimToRealGapLCOnly{}$ for LC-only
($\TierLCOnlyMacroFFull{} \to \LCOnlyMacroF{}$). It is read here as a
result, not a failure: Tier A carries none of Tier B's instrument noise,
real-bogus artifacts, or host-catalogue error, so a model that scored the
same on both would be the stronger evidence of a broken, leaking evaluation
split, since nothing about Tier B should be \emph{easier} than its own
simulated approximation. Part of the gap is a plain brightness mismatch:
real BTS objects are far brighter than the simulated training population
(Section~\ref{sec:limitations}). The transfer gains over the matching
zero-shot models ($+\TransferGain{}$ for LC+static, $+\TransferGainLCOnly{}$
for LC-only) show that fine-tuning is not cosmetic: a model carrying only
Tier A weights, evaluated on Tier B with no adaptation at all, is
measurably worse than one given even a small amount of real labelled data.

\paragraph{Ib/c misses its pre-registered band.} Before any Tier B result
existed, this project expected Ib/c F1 between \IbcBandLow{} and
\IbcBandHigh{}. The deployed configuration reaches \WithStaticFIbc{} (mean
over seeds; \DeployedFIbc{} for the served checkpoint), and LC-only reaches
\LCOnlyFIbc{}: both are below the band. The band is reported unchanged and
the miss is reported as a miss. Two reasons are likely, neither of which
would justify moving the band after the fact. First, the Tier B test split
holds only \TierBSupportIbc{} Ib/c objects against \TierBSupportIa{} Ia, so a
handful of objects moves this F1 substantially. Second, the band was
derived from published models trained and evaluated with far more Ib/c
examples than the \TierBJoinCountIbc{} spectroscopically confirmed Ib/c
objects available here across all splits. Before the zeropoint correction
this figure was \PreFixWithStaticFIbc{} for LC+static and
\PreFixLCOnlyFIbc{} for LC-only, inside the band, but those values came from
mis-scaled input and are not evidence that the band was met.

\section{Host features help on both tiers}

With correctly scaled fluxes, host and static features improve
classification on both tiers: by $\PhotozGainWith{}$ macro-F1 on Tier A
($\TierLCOnlyMacroFFull{} \to \TierLCStaticMacroFFull{}$) and by
$\StaticGainTierB{}$ on Tier B ($\LCOnlyMacroF{} \to \WithStaticMacroF{}$,
\StaticGainOverSD{} times the larger seed-to-seed standard deviation).
LC+static is therefore the deployed configuration.

An earlier version of this analysis reported the opposite on Tier B -- host
features making real-data fine-tuning worse -- and treated it as an open
finding. That inversion was an artifact of the flux-zeropoint conversion
error described in Section~\ref{sec:silentfailures}. Under identical code,
seeds and object order, the old fluxes put LC-only
\OldFluxStaticInversion{} macro-F1 above LC+static
($\OldFluxLCOnlyMacroF{}$ against $\OldFluxWithStaticMacroF{}$); the
corrected fluxes reverse the order. There is no remaining inversion to
explain.

Separately, on Tier A only, retraining LC+static without host photometric
redshift keeps \PhotozGainFractionSurviving{}\% of the static-feature gain
($+\PhotozGainWithout{}$ macro-F1 without photo-$z$, against
$+\PhotozGainWith{}$ with it): photo-$z$ contributes, but most of the Tier A
benefit comes from the other host fields.

\section{Windowing: one mechanism, three symptoms}

The train/serve sequence-length mismatch (Section~\ref{sec:windowing}), the
unbounded right edge for non-fading sources (Table~\ref{tab:spans}), and
classifier entropy's inversion under domain shift are one mechanism
observed at three points, not three separate findings. A model trained on
windowed sequences averaging \WindowMeanSeqLenWindowed{} observations
(Figure~\ref{fig:windowhist}) meets, at serving time, a changing-look AGN
whose uncapped tensorized sequence runs to a mean of \ClagnSeqLenUncapped{}
observations -- an input far outside anything the encoder was ever fit
against. The measured consequence is not increased uncertainty but
increased, misplaced confidence: CLAGN's mean entropy under the uncapped
window (\ClagnEntropyUncapped{}) sits \emph{below} the Tier A in-distribution
test population's own baseline (\EntropyTierANormal{}), and capping the
window at \SpanChosenDefault{} days moves it substantially toward, though not all the way
to, that baseline (\ClagnEntropyCapped{}). The general lesson is that a live
serving pipeline needs a bounded, \emph{recent} window for any long-lived
source, not a window defined relative to first detection alone -- a
distinction this project's own span cap addresses for training-distribution
matching but does not, by itself, solve for a serving-time rolling window
(Section~\ref{sec:analysis}).

\section{Confidently wrong on out-of-distribution input}

Table~\ref{tab:entropybyclass} makes the general pattern explicit across
every held-out simulated exotic class: mean entropy is markedly lower --
i.e., the classifier is more confident -- on several classes it was never
trained to recognise than on its own in-distribution test population
(\EntropyTierANormal{}). This is precisely the failure mode an anomaly
detector exists to catch, and precisely why classifier confidence cannot be
the component that catches it: a component that becomes \emph{more}
confident, not less, on the population it is meant to flag will actively
mislead a combination that trusts it, which is exactly what dropping it
measurably improved (Table~\ref{tab:anomaly}).

\begin{table}[htb]
\centering
\caption{Mean classifier entropy per held-out Tier A exotic class (computed
directly from the per-object entropy values already scored for the anomaly
branch, grouped by taxonomy id), against the Tier A in-distribution test
population's own baseline. CLAGN, the class with by far the longest
detection span (Table~\ref{tab:spans}), is also the most confidently wrong;
KN, a genuinely fast transient despite being outside the 3-class target
set, shows no inversion at all.}
\label{tab:entropybyclass}
\begin{tabular}{lrr}
\toprule
Population & $N$ & Mean entropy \\
\midrule
Tier A test (in-distribution) & \AnomalyTierANNormal{} & \EntropyTierANormal{} \\
KN & \ExoticNKN{} & \ExoticEntropyKN{} \\
CART & \ExoticNCART{} & \ExoticEntropyCART{} \\
SLSN & \ExoticNSLSN{} & \ExoticEntropySLSN{} \\
uLens & \ExoticNuLens{} & \ExoticEntropyuLens{} \\
ILOT & \ExoticNILOT{} & \ExoticEntropyILOT{} \\
TDE & \ExoticNTDE{} & \ExoticEntropyTDE{} \\
PISN & \ExoticNPISN{} & \ExoticEntropyPISN{} \\
CLAGN & \ExoticNCLAGN{} & \ExoticEntropyCLAGN{} \\
\bottomrule
\end{tabular}
\end{table}

\section{Anomaly components under domain shift}

The two latent-space methods (isolation forest, Mahalanobis distance) do
not survive the shift from simulated to real data because the latent space
they operate on is the Tier A encoder's own -- and that geometry simply
does not transfer to real ZTF objects the encoder never saw a shred of
during training. Reconstruction error moves the other way, because it
measures something closer to the physical light curve itself: a
cataclysmic variable or nova genuinely does not resemble a supernova's
flux-versus-time behaviour, a property of the raw signal that survives
exactly the domain shift that breaks the two latent-geometry methods. The
combination should be read as \emph{robustness under domain shift}, not as
peak performance -- it is beaten by the single best component on each
individual set (isolation forest on Tier A, reconstruction error on Tier B)
and is the only score that stays useful on both: its weaker value
(\AnomalyTierBCombinedThreeRecomputed{}, Tier B) is higher than the weaker of
the two values for any single component.
That is the honest, and more useful, claim to make about it.

\section{Systems findings}

The classifier itself is a small fraction of what an end-to-end serving
pipeline actually spends time on: ONNX inference takes
\OnnxMsPerAlertBatchOne{}\,ms per alert against a median end-to-end latency
of \ServingLatencyPFifty{}\,ms, which is dominated by the per-alert Redis
and Postgres round trips in this implementation
(Section~\ref{sec:results}). The implication for a broker-scale deployment is that
batching the state-layer writes, not further optimising the model, is the
obvious next lever if higher throughput is ever required.

\section{Validation findings: what would have failed silently}
\label{sec:silentfailures}

Table~\ref{tab:silentfailures} is, in its way, this report's methodological
contribution: every item in it would have produced a plausible, confident,
\emph{wrong} result, and none of them would have raised an error or crashed
anything.

The first row is qualitatively different from the rest, and it is the
strongest example of this section's argument. The project's governing
document described the conversion from a survey's native zeropoint to the
common FLUXCAL scale as a multiplication by a fixed factor, when the
correct operation is a division. That wrong value propagated into the
implementation, and then into the unit test written to verify the
implementation, which asserted the same factor. Implementation and test
agreed with each other and with the document, so all \TestsPassedPreFix{}
tests passed while every real ZTF flux was stored a median
\FluxBugFactor{} times too faint (\FluxBugFactorLow{}--\FluxBugFactorHigh{}
times across rows, since the ZTF zeropoint varies per exposure) and every
live LSST flux \LsstInversionFactor{} times too bright. A golden test is
only as independent as the number it was computed from; this one was
computed from the same wrong statement as the code. The replacement test
checks a property rather than restating a formula: the AB magnitude of a
source must be unchanged by the conversion. Once it existed,
\TestsFailedInvarianceOnOldCode{} tests failed against the old code. The
error survived every pre-registered gate because each Tier B result was
internally consistent -- fine-tuning, testing and serving all saw the same
mis-scaled fluxes -- and it surfaced only in comparisons across tiers,
where it produced a spurious finding (Section~\ref{sec:discussion}, host
features) and the wrong choice of deployed model.

\begin{table}[htb]
\centering
\caption{Silent-failure modes identified during development: what each was,
how it was caught, and what would have happened unchecked. None would have
raised an error.}
\label{tab:silentfailures}
\scriptsize
\begin{tabular}{L{2.7cm}L{4.6cm}L{5.8cm}}
\toprule
Issue & How it was caught & Consequence if unchecked \\
\midrule
Zeropoint conversion inverted in the governing document, the code and its
  test & A post-hoc audit checking AB-magnitude invariance numerically,
  after all \TestsPassedPreFix{} tests had passed & Real-data fluxes
  mis-scaled by an order of magnitude; a spurious ``host features hurt''
  finding and the wrong deployed model \\
Undocumented sentinel conventions (Table~\ref{tab:sentinels}) & Manual
  cross-check of raw FITS values against the format spec & Missing host
  features would train as real, extremely faint measurements \\
Three coexisting numbering schemes & Direct inspection found three
  different values for one object & Training against scrambled,
  meaningless class labels, indistinguishable from ``the model isn't
  learning'' \\
Full-history training data vs.\ a windowed live stream &
  Direct measurement of mean detections vs.\ forced photometry
  (Section~\ref{sec:windowing}) & A systematic train/serve distribution
  mismatch invisible in any offline metric, since the offline test split
  shares the training set's own (wrong) full-history shape \\
Unapplied/assumed-constant flux zeropoint & Per-row \texttt{magzpsci}
  variation measured directly on real data, not assumed constant &
  Silent, brightness-dependent flux error, worst in the faint early-time
  regime this project prioritises \\
A binary triage flag mistaken for a subtype label & Direct inspection of
  BTSbot metadata's column semantics before use & Training against ``is
  this a real transient at all,'' not the intended 3-class question,
  producing plausible-looking numbers for the wrong task \\
Inverted entropy component & Direct comparison of mean entropy,
  in-distribution vs.\ held-out anomalous, both tiers & A published
  anomaly score quietly worse than it needed to be, with no failure
  symptom pointing at entropy specifically \\
Undersized CCD-edge cutouts discarding an entire alert & A live ingestion
  run (Section~\ref{sec:results}), not offline testing & Live alerts'
  light curves discarded over an unrelated cutout-shape issue \\
\bottomrule
\end{tabular}
\end{table}

The common property across every row is that none of them would have raised
an error; each would have produced a confident, plausible, wrong result.
What caught them was not a single tool but reinforcing habits: a frozen
evaluation harness that could not be quietly adjusted to make a
disappointing number look better (Section~\ref{sec:methodology}), golden
tests written against one hand-verified real object before the reader code
that would have to pass them existed, pre-registered gates
(Table~\ref{tab:gates}) defined before the results they judge, and -- the
lesson of the first row -- tests that check a physical invariant rather
than a transcribed constant. This is the project's actual methodological
contribution, and is presented as such rather than as an apology.

\section{Limitations}
\label{sec:limitations}

Real BTS objects are much brighter than the simulated training set: the
median detection flux is \FluxMedianTierB{} FLUXCAL on Tier B against
\FluxMedianTierA{} on Tier A, a factor of about \FluxBrightnessRatio{}. The
conversion error masked roughly \FluxBugFactor{} times of that difference
(the pre-fix Tier B median was \FluxMedianTierBPreFix{}), so the encoder
was fine-tuned with a frozen-trunk, low-learning-rate schedule chosen
before the true size of the brightness shift was visible. This mismatch is
a plausible contributor to the sim-to-real gap; it has not been isolated,
and no schedule designed for it was tried.

Tier B is a systematically cleaner slice of the real sky than a live
stream, not a representative sample of one: it is pre-filtered to
brightening subtractions before a candidate ever reaches its public
release (\TierBNegFluxFrac{}\% negative flux, against
\LiveNovZtfNegFluxFrac{}\% measured on live ZTF traffic), it contains no
forced-photometry row type at all, and its labels are drawn from a
magnitude-limited spectroscopic follow-up programme, not from LSST's own
future detection limit.

Output probabilities carry the deployed model's training prior, not the
sky prior. On \LiveNovZtfNObjects{} live ZTF objects the deployed model
predicts Ia \LiveNovZtfRawIa{}\%, Ib/c \LiveNovZtfRawIbc{}\% and II
\LiveNovZtfRawII{}\%; Ib/c sits below its real sky rate, consistent with the
low Ib/c recall above. A single multiplicative prior correction does not
repair this: dividing out the ELAsTiCC prior gives Ib/c
\LiveNovZtfElasticcIbcRangeLow{}--\LiveNovZtfElasticcIbcRangeHigh{}\%
across plausible sky priors, and dividing out the effective Tier B
fine-tuning prior (Ia \TierBEffectivePriorIa{}\%, Ib/c
\TierBEffectivePriorIbc{}\%, II \TierBEffectivePriorII{}\%) pushes II to
\LiveNovZtfTierBCorrII{}\%. Calibrating the output probabilities is a
prerequisite for any rate estimate and has not been done.

LSST ingestion now works on real traffic (Section~\ref{sec:results}), but
LSST alerts carry no labels, so no accuracy claim is possible for them,
and their cutout stamps (\LsstArchiveStampMin{}--\LsstArchiveStampMax{}
pixels on a side) do not fit the fixed cutout shape of the canonical
schema, so LSST cutouts are not stored. The image encoder was cut, not
deferred: Tier B image data was never loaded across the project, and no
schedule remained to build, fuse, and retrain around it once that became
clear. Gate D remains unresolved, for the reason given in
Table~\ref{tab:gateoutcomes}.
